\section{怎样修改文章}

古今中外，凡是文章写得好的人，恐怕没有不在修改上下过很大功夫的。说某作家文不加点，一挥而就，除了强调其人文思泉涌，下笔敏捷之外，大约都是由于在打腹稿时不断修正的结果。比如号称“初唐四杰”之首的王勃，要写文章时，总是仰面朝天地躺在床上，扯条被子连头带脚蒙上，好象睡着了，其实在紧张、亢奋地构思、推敲文章呢！过了一阵子，他猛地掀掉被子，跳下床来，就能手不停挥地写好一篇文章，甚至不用改动一个字。一部《红楼梦》，人们都说写得好，殊不知作者曹雪芹呕心沥血地写了十年，中途还修改了五次呢。据说俄国大文学家列夫·托尔斯泰在创作小说《复活》时，对女主人公玛丝洛娃在法院里的形象的描写刻画就修改了二十次之多；法国十九世纪的伟大作家巴尔扎克的每部作品都要改六、七次，而且在排印时他还要修改，有时修改的部分比原来的本文还要多，甚至每部小说再版时，他还要重新修改，连细节也不肯放过。对于我们初学写作者来说，学会自己动手修改文章就更为重要了，因为我们本来就没有那些大作家的高水平，文章毛病多，而且也只有经过自己的一番深入思考、反复琢磨，才可以逐渐明了为什么这样写好，那样写不好，逐渐提高自己的写作水平。

那么，我们怎样着手修改文章呢？当然，第一步是“找毛病”，第二步是“动手术”。具体说来，主要应注意以下几点：

\textbf{1. 推敲文章的主题是否正确、深刻、新颖，材料是否用得恰当、充实、典型。}

如果说象王勃那样蒙头大睡打腹稿是一种独特的构思、推敲文章的方式，我们则可以根据作文提纲来琢磨文章的立意和材料的运用。这种涉及全局的修改，必须在下笔之先就该做好；到了作文写成之后再来修改，时间可就来不及了。所谓“文贵立意”，“意在笔先”，正是说明这个问题的重要。因此，我们每次作文时，必须要深入地分析材料，准确地把握材料的实质，从中选取典型的材料来提炼主题，深化主题，如果发现所用的材料不能正确地表现主题，或者只能一般化地表现主题，那就该毫不留情地删去，或者修正原来材料的角度，或者补充更能表现主题的材料，务求能使主题不仅正确，而且有较为深刻的思想意义，较为新颖的思想角度。例如，有位同学在写《我们的班主任》这篇作文时，他先拟好了这样一份作文提纲：

（一）给班主任画像，说明他是一个对自己刻苦、对同学严厉的人。

（二）他上课时，谁做个小动作，都要挨批评，可一考试，好多题目总象是被他猜中了似的。我班学习成绩经常居全年级第一。

（三）每次运动会前，他都带着我们锻炼，总分第一名的锦旗总由我班保持。有一次，陈华因点名未到被取消了赛跑资格，他把陈华都批评哭了。

（四）有次劳动时，他只吃同学们剩下的馒头，还批评同学们浪费。

（五）结束语（略写）。

他拟好了提纲之后，没有急着动笔，而是对提纲进行推敲和修改。他发现这些材料虽然都是实事，但是能表达一个什么样主题呢？这样的主题又是否正确、深刻、新颖呢？他觉得有问题，关键在于只抓了现象，没抓住实质。于是他深入地分析材料，不断地开掘下去，认识到这位班主任实在是一位严格要求同学们在德、智、体诸方面全面发展的好“园丁”。立意提高之后，他又对每项材料换取角度进行表述，并且围绕这个新的主题补充了一些新材料。于是，他选取“严格的爱”这个角度，深刻而又新颖地表达了自己的构思：

（一）我班事事走在前面的谜——谜底就是有这样一位对我们“严格的爱”的班主任。

（二）“课堂就是战场”——班主任的口号和我现。

（三）“枪声响起的时候，士兵怎能没有在战斗岗位上？”——班主任对陈华的疏忽进行批评，使陈华惭愧得哭了起来，奋勇夺取长跑第一名。

（四）“谁知盘中餐，粒粒皆辛苦”——班主任吃同学们剩下的馒头，身体力行，使同学们勤俭成风。

（五）“骄傲使人落后”——当我获得物理竞赛第一名感到轻飘飘时，班主任及时敲响了警钟。

（六）以抒情笔调结尾，表达对这位好“园丁”的敬爱心情。

请看，这两份提纲几乎用同样的材料，但是由于所选取的角度不同，所开掘的程度不同，二者立意的高度和深度就可不同了。这种类型的修改牵动了全文，可以叫做“大改”吧。“大改”一般在构思、编写提纲的过程中进行。列宁的夫人克鲁普斯卡娅曾回忆说：“他（列宁）在写文章之前，通常是先写好提纲。从这个提纲可以追溯伊里奇的整个思想过程。有许多文章，伊里奇把提纲改了两遍、三遍。”列宁的做法，不正可以给我们以极大的启示吗？写文章时，中途也往往会发生构思的改变，如果文章写出一半，发现问题去改，就费事了。这象房屋建了一半，推倒重来就麻烦；而在蓝图上改动，那就轻而易举了。
\newpage

\textbf{2. 调整文章的结构，修改文章中不够连贯、缺乏条理、前后颠倒、详略不当等有毛病的部分。}

文章的结构和层次，本来在编写作文提纲时已经体现了出来，如有不妥之处，可以在修改提纲时予以修正。不过，提纲毕竟只是理清思路、组织材料、安排层次的蓝图。蓝图尽管是未来房屋的模型，还不是真正的房屋。哪儿该接榫，哪儿用铆钉，哪儿该强调装饰，哪儿该力求朴素，先立哪根柱，后放哪根梁，等等，这在蓝图中是没法体现出来的。所以，我们在拟定提纲动笔作文的过程中，往往会发现某些缺之

过渡或过渡不自然，缺乏必要的照应，开头不切题或太突兀，结尾一般化或拖沓无力，有时还收不到尾；甚至会发现该详的没详，该略的没略，该先说的放到了后面，该后说的放到了前面，该分段的没分，不该分段的又分了，该倒叙的用了顺叙，该顺叙的又用了倒叙，倒叙、插叙的起止界限不分明，划分自然段太零碎或者太笼统，自然段有的过长，有的过短，显得不够匀称，不够完整，等等。这些毛病，尽管还不及构思立意那样要牵动全文，但是由于层次不清，条理不明，主次不分，详略不当，缺乏过渡和照应，开头不象“风头”那样漂亮、俊秀，主干不象“猪肚”那样饱满、充实，结语不象“豹尾”那样结实有力，就影响到全文的结构，达不到严谨自然，完整统一的高度，从而就不能很好地表达中心思想，所以必须认真地琢磨，下细的推敲，一一予以修改，直到改得满意才罢手。

原高中语文课本第三册上，有何其芳同志的《谈修改文章》。文中列举了十二种毛病，有的属于内容方面的（主要是主题和材料），有的属于形式方面的（主要是结构和语言），可以供我们“找毛病”时参考。由于结构上的毛病，既不如构思立意那样范围大，牵动全文，又不象斟酌词句那样范围小，限于一字一句，所以我们可以把关于文章结构的修改叫做“中改”吧。这种修改，不仅在编写作文提纲时就该着手了，而且在写作过程中需要不断地进行，一直坚持到作文的完成。

\textbf{3. 修改文章中不合逻辑、不合语法、修辞不当的地方。}

毛泽东同志说：“总之，一个合逻辑，一个合文法，一个较好的修辞。这三点请你们在写文章的时候注意。”毛主席要我们在写文章的时候应该注意的三点，正是一般所说的修改文章的要点，对于我们初学写作者来说，尤其显得重要。

合乎逻辑，是说文章的开头、中间、结尾不能互相矛盾，要有紧密的内在联系，使全文形成统一的整体；也包括文章中的一句话、一段话前后一致，合乎事理。例如，有的作文中，前面写别人称呼“我”为“小李”，后面写别人称呼“我”为“小王”，这就使读者莫名其妙了。再如，前面写张明早早地到校，拿着提桶去接水，准备打扫教室的卫生；接着写班长到校，自言自语地说：“奇怪，是谁把提桶拿走了？”后面却写“原来张明头晚就与班长约好，今早一起把教室打扫一下。”如果是“头晚约好”的，班长又怎么感到“奇怪”呢？又如，有篇作文前面写的是要发扬艰苦奋斗的作风，后面举的事例却是个人注意锻炼身体，从不生病，节省了医药费。不生病，节省药费，与艰苦奋斗没有紧密的内在联系，材料说明不了观点。这类的写法，都不能使文意前后圆合，甚至互相矛盾，当然就不合逻辑了。对于这类不合逻辑的地方，我们要统观上下文，选准统一的角度，把它们修改得前后一致才行。至于一句话、一段话写得不合逻辑的现象就更普遍了。例如：有篇作文，描写丰收的景象说：“火红的高粱，金黄的麦子，迎风摇曳，点头微笑。”高粱晒红米的时候，麦子早已割完了，两者不能同时“迎风摇曳”。这段描写，本身矛盾，不合实际事理，也就是不合逻辑。再如：有篇作文，描写暴风雨前景象说：“在那一碧如洗的蓝天上，飘浮着朵朵白云。突然，一阵风过，乌云翻滚，天昏地暗，一场大雨即将来临了。”既然是“一碧如洗的蓝天”，怎能有“朵朵白云”？“朵朵白云”又怎能“突然”间成了“翻滚”的“乌云”？这种前言不搭后语的描写，违反人们的常识，也是不合逻辑的。关于这种不合逻辑的病句，我们在《怎样修改病句》中已经作了怎样修改的说明，重要的是应把已经掌握的修改这类病句的方法，灵活地运用到修改作文当中去。

合乎文法，是指用词准确，词语搭配得当，句子结构完整。对于初学写作者来说，写通句子是写好文章的第一步。很难想象，病句比比皆是却能够写好文章。在有些同学的作文中，不是经常出现“通过参观法制教育展览会，给我上了一堂生动的政治课”或者“校园的春天是一个多么理想的学习园地啊”这样一类缺少句子基本成分或者前后配搭不当的病句吗？我们在写作文时，应该一词一句地进行推敲，决不可马马虎虎，写完了事，要做到字斟句酌，精益求精。我想，只要我们态度认真，有意识地运用已经掌握的修改病句的知识和方法，联系上下文意，分析句子结构和词语之间的配搭，自能发现句子的毛病，并且能够加以改正。例如有这样一篇广播稿《向郭瑾学习》（见302页），我们只要用心琢磨，就能指出文中的毛病，说明自己的理由，并加以改正的。（见下页）
\newpage

至于较好的修辞，主要是指语言运用得准确、鲜明、生动。毛泽东同志曾举过一个典型的例子，说：“一九五六年国家预算报告中说过‘稳妥可靠’这个话，我建议以后改为‘充分可靠’……稳妥和可靠，意思是重复的。用‘稳妥’形容‘可靠’，没有增加什么，也没有限制什么。形容词一面是修饰词，一面是限制词。说‘充分可靠’，这就在程度上限制了它，不是普通可靠，是‘充分可靠’。”语言学家把这种不用修辞格的修辞叫做消极修辞，把用修辞格的修辞叫做积极修辞。其实消极修辞比积极修辞用得更为广泛，我们一定要努力在这方面多下功夫，把文章改得文从字顺，语言准确、鲜明、生动。对于积极修辞，我们当然也不能放过，但是一定要注意运用得当，细心推敲那些比喻不当、比拟不伦、夸张失实等语病，并认真改正。

\noindent
\begin{tabular}{|p{12em}|p{12em}|}
	\hline
	原稿 & 评改 \\ \hline
	郭璞同学人小志气大。她不管体质软弱，但是天天第一个到校。她说：“今天的刻苦学习，是为了未来参加四个现代化的建设。我们精力充分，怎能不勤奋的学习呢？” & “身小”仅指个子矮；不如改用“人小”，兼指岁数和个子。“不管”与“但是”配搭不当。“不管”应改为“尽管”。这里是坚持不懈进行下去的意思，应该改用“坚持”。“的”字多余，删去。“未来”改为“明天”。“明天”与“今天”对称，又有紧迫感。“充分”改为“充沛”。“充沛”是充足、旺盛的意思，形容精力合适。“的”字多余，删去。\\ \hline
	她说到做到，各门功课都学得很好，难怪老师们没有一个不夸她不是一个好学生。现在，郭璞这种刻苦学习的精神已成了我们学习的好榜样，在她的带动下，我班为四化而勤奋学习已蔚然成风了。 & “没有”、“不”、“不是”，多重否定，造成混乱，意思全反。应将“不是”改为“真是”。“精神”与“好榜样”配搭不当。应该改为“现在，郭璞已成了我们学习的好榜样。在她刻苦学习的精神带动下，我班······”。 \\ \hline
\end{tabular}
\vspace{1em}

为了启发同学们怎样给自己的作文“找毛病”和“动手术”，我们将著名语言学家朱德熙先生评改的一段中学生

作文引述如下，请同学们细心体会：

[\ 原文\ ]

成熟的西红柿鲜嫩可口，营养丰富。人们爱吃它，细菌也爱吃它。正
\uline{因为它营养丰富，水分又大}①，细菌
\uline{一旦}②在它里面安家落户，传宗接代，它立刻就要腐烂。所以，贮存的
\uline{方法}③，第一是灭菌，第二是
\uline{使细菌无法再行侵入}④，
\uline{不能再繁殖}⑤。
\uline{为了灭菌}⑥，家庭普通的蒸锅就够了。把西红柿切碎，装入瓶内，上锅一蒸，就可以将细菌杀死。但是，如果蒸的时间太长，西红柿就不能保持
\uline{原形}⑦，
\uline{失去了外形的美观}⑧；蒸的时间太短，又不能保证将细菌全部消灭，还有腐烂的可能。因此，
\uline{以蒸十五分钟到三十分\\ 钟为宜}⑨。
\begin{flushright}
	（北大附中孙思东《万花纷谢柿犹红》片断）
\end{flushright}


[\ 评改\ ]

①\ “营养丰富”四字与上文重复，而且说西红柿容易腐烂，是因为它“营养丰富”，也不合适。这十二个字不如删去。

②\ “一旦”挪至句首。“安家落户”、“传宗接代”，总需要些时间，说“立刻就要腐烂”不妥，“立刻”改为“很快”。

③\ “方法”改为“要点”。

④\ ⑤\ 这里是说的贮存方法，而“使细菌无法再行侵入”并不是一种方法，不如改为“第二是密封，使空气里的细菌无法侵入”。此外，“无法再行侵入”和“不能再繁殖”两件事毫无关系。尽管细菌不能再“侵入”瓶内，它照样可以在瓶子外边繁殖。可见“不能再繁殖”一句是多余的，应删去。
\newpage

⑥\ “为了灭菌”改为“灭菌并不需要什么特殊的设备”。

⑦\ ⑧\ 上文说“把西红柿切碎，装入瓶内”，既然已经切碎了，怎么还能保持原形呢？“原形”改为“原来的色泽”，删去“失去了外形的美观”一句。⑨蒸的时间不能太长，也不能太短，从这一点并不能得出“以蒸十五分钟到三十分钟为宜”的结论。原文用了“因此”，就显得前后两句话似乎有因果关系，其实蒸多长时间合适，完全是从实践中得来的经验。这一句可以改为：“因此蒸的时间要掌握好，一般以十五分钟到三十分钟为宜。”

我们看，朱先生对这段文字的语病推敲得多么细致、多么入理啊！尽管比起构思立意的“大改”，结构条理的“中改”来，这种字斟句酌甚至连一个标点符号也不放过的修改只是“小改”，但是它在作文训练中又是多么重要的基本功啊。人们所说的狭义的修改文章指的正是这种类型的“小改”。我们如果能独立改好自己作文中存在的语病，那就说明你已具备比较扎实的写作基本功了。
\newpage